\section{Introducción}\label{sec:intro_informe_final}

El laboratorio del Departamento de Ingeniería en Computación (LabDIC) de la Universidad de Magallanes gestiona actualmente los préstamos de material tecnológico mediante formularios impresos y registros manuales, lo que dificulta el control, genera duplicidad de información y limita la trazabilidad. Ante esta situación, el presente proyecto propone el desarrollo de una plataforma web que digitalice estos procesos, mantenga actualizada la base de datos del inventario y permita a encargados y solicitantes gestionar las solicitudes de préstamo de forma eficiente. El primer paso consistió en establecer las bases conceptuales, tecnológicas y de planificación que orientan el desarrollo, incluyendo la selección del stack tecnológico, el análisis de requerimientos funcionales y no funcionales, y la estructuración del trabajo en hitos progresivos.

El diseño del modelo de datos partió de tres grandes ejes funcionales ---usuarios, productos y préstamos--- que, al profundizar en las interacciones entre entidades, derivaron en seis módulos: Usuarios, Autenticación, Catálogo, Productos, Dispositivos y Préstamos. Este diseño modular, implementado con SQLAlchemy como ORM y Alembic para el control de migraciones, garantiza que el esquema pueda evolucionar de forma documentada y reproducible a lo largo de todo el ciclo de desarrollo.

Antes de abordar la implementación de cada módulo, se documentan las decisiones arquitectónicas que estructuran el sistema. En el backend, se define la separación de responsabilidades entre DTOs, repositorios y controladores agrupados bajo el concepto de \textit{servicio}; en el frontend, se documentan las decisiones equivalentes: Vue Router para navegación y control de acceso, Pinia para estado global persistente, y las carpetas \texttt{types/} y \texttt{services/} como puente tipado hacia la API. El flujo de consumo de la API sintetiza visualmente cómo los datos recorren el sistema completo, desde PostgreSQL hasta los componentes de la interfaz.

La implementación del backend desarrolló los servicios, repositorios, DTOs y controladores para los seis módulos funcionales, exponiendo una API REST completa y documentada. El proceso no fue lineal: a medida que cada módulo tomaba forma, surgieron ajustes en el modelo de datos gestionados mediante migraciones con Alembic, evidenciando la naturaleza iterativa del desarrollo.

El frontend se construyó en seis fases progresivas que van desde la autenticación y control de acceso hasta la gestión completa de préstamos y la revisión de UX y UI. Durante esta etapa se confirmó la interdependencia real entre ambas capas del sistema: la implementación de vistas de productos y dispositivos reveló la necesidad de correcciones en cascada en la base de datos. Finalmente, se documenta el trabajo futuro, identificando mejoras pendientes al sistema actual y proponiendo nuevos módulos que la arquitectura modular de la plataforma está preparada para recibir.